\chapter{Sharp Runtime: Strings, Streams, and a Real Audit History}
\label{ch:sharp-runtime-namespaces}

This chapter covers two smaller, practical corners of \texttt{sharp-runtime}'s API --- its
text and stream types --- and then the part of its history most worth learning from: what
happened when the project's own completed task tracker was deliberately \emph{not} trusted as
proof of correctness, and independent audits went looking for what it might have missed.

\section{String and StringBuilder: a static-only utility versus a real mutable buffer}

\cnaclass{System::String} is deliberately \emph{not} a wrapper class at all --- its default
constructor and destructor are both explicitly deleted, making it a pure static-method utility
operating on ordinary \texttt{std::string} values (\texttt{IsNullOrEmpty}, \texttt{Compare},
the \texttt{IndexOf} / \texttt{LastIndexOf} family, \texttt{Split}, \texttt{Substring},
\texttt{Trim*}, \texttt{PadLeft} / \texttt{PadRight}). \cnaclass{Text::StringBuilder}, by
contrast, is a real, stateful, mutable buffer type --- the many-overloaded \texttt{Append}
(string, char, int, double, float, bool, and the \texttt{longcs} alias from
Chapter~\ref{ch:design-philosophy}), \texttt{AppendLine}, \texttt{Insert},
\texttt{Remove}, \texttt{Replace}, and a nested \texttt{ChunkEnumerator} class mirroring real
.NET's own chunked-buffer enumeration shape --- real .NET's own \texttt{StringBuilder} is
internally a linked list of separate character chunks, and \texttt{ChunkEnumerator} walks that
list without copying; this port's \texttt{StringBuilder} is backed by a single
\texttt{std::string} instead, so its own \texttt{ChunkEnumerator} always yields exactly one
chunk holding the entire current buffer --- the enumeration \emph{shape} is preserved
faithfully even though the simpler single-buffer implementation underneath never actually
produces more than one chunk to walk.

\begin{lstlisting}
// Worked example: System::String is a pure static utility -- there is no
// instance to construct, only free functions operating on std::string.
std::vector<std::string> parts = System::String::Split(
    "Content/Textures/hero.png", '/');
std::string trimmed = System::String::Trim(parts.back());

// Text::StringBuilder, by contrast, is a real mutable buffer.
System::Text::StringBuilder sb;
sb.Append("Score: ").Append(playerScore).AppendLine();
sb.Append("Lives: ").Append(livesRemaining);

for (const std::string& chunk : sb.GetChunks())
{
    // Exactly one iteration, always -- see the ChunkEnumerator note above.
    RenderHudText(chunk);
}
\end{lstlisting}

\section{Stream: a lightweight subset, honest about it}

\cnaclass{IO::Stream}'s own class comment states its scope plainly: ``lightweight subset of
the .NET Stream API.'' \texttt{Read(buffer, offset, count)}, \texttt{Close()}, and
\texttt{getLengthProperty()} are pure virtual; \texttt{Write} / \texttt{WriteByte} default to
throwing \texttt{NotSupportedException} unless a writable stream subclass overrides them ---
the same shape real .NET's own read-only stream implementations follow, rather than forcing
every stream subclass to implement a full read/write/seek contract regardless of whether it
actually supports all three. \texttt{getCanWriteProperty()} defaults to false and
\texttt{getCanReadProperty()} defaults to true, matching a read-only stream's honest self-report
without requiring every subclass to override both just to state the obvious.

\begin{lstlisting}
// Worked example: the minimal real subclass this base class actually
// requires -- only three overrides, matching a genuinely read-only
// stream's real capabilities (Write/WriteByte are correctly inherited as
// NotSupportedException-throwing, never called by a well-behaved reader).
class MemoryBlobStream final : public System::IO::Stream
{
public:
    explicit MemoryBlobStream(std::vector<SharpRuntime::bytecs> data)
        : data_(std::move(data)) {}

    SharpRuntime::intcs Read(SharpRuntime::bytecs buffer[],
                             SharpRuntime::intcs offset, SharpRuntime::intcs count) override
    {
        SharpRuntime::intcs n = std::min<SharpRuntime::intcs>(
            count, static_cast<SharpRuntime::intcs>(data_.size()) - position_);
        std::memcpy(buffer + offset, data_.data() + position_, n);
        position_ += n;
        return n;
    }

    void Close() override {}

    [[nodiscard]] SharpRuntime::intcs getLengthProperty() const override
    {
        return static_cast<SharpRuntime::intcs>(data_.size());
    }

private:
    std::vector<SharpRuntime::bytecs> data_;
    SharpRuntime::intcs position_ = 0;
};
\end{lstlisting}

\section{The post-stabilization audit: distrusting a complete-looking task list on purpose}

\texttt{sharp-runtime}'s own ticket backlog --- 1{,}709 tickets across sixteen categories ---
reached 100\% \texttt{done} at one point in the project's history. What happened next is the
single most instructive episode in this project's own story, and worth recounting in full,
because it is the same discipline Chapter~\ref{ch:project-practice} named as this whole
ecosystem's most distinctive trait, applied at real scale: rather than treating a fully
checked-off ticket table as proof the runtime was correct, the project commissioned a fresh,
independent audit specifically built to distrust that conclusion. Seven parallel, read-only
audit agents were dispatched across eight risk categories --- API inconsistencies,
undocumented deviations from real .NET, remaining \texttt{NotImplementedException}s, partial
or stub classes, platform-specific issues, exception-type mismatches, silently-wrong
behavior, and missing high-risk test coverage --- each one required to verify every finding
directly against real .NET source and against the project's own documented, accepted
deviations list before it could be reported at all, specifically to filter out
plausible-sounding but wrong findings before they ever reached a human.

The result was twenty concrete findings --- nine high-severity, eight medium, three low ---
filed as new tickets and subsequently fixed. One is worth naming in full, because of how
quietly it could have shipped: \texttt{Convert::ToXxx(double/float)} was implemented as a
bare \texttt{static\_cast<T>()}, meaning simple truncation rather than round-half-to-even
rounding. A direct repro confirmed \texttt{Convert::ToInt32(2.9)} returned \texttt{2} where
real .NET returns \texttt{3}, and \texttt{Convert::ToInt32(3.5)} returned \texttt{3} where real
.NET returns \texttt{4} --- and, more strikingly, \texttt{ToUInt32(double)}'s own doc comment
already said \texttt{"(truncates)"}, as though the wrong behavior had been documented as
intentional rather than caught as a bug.

\begin{lstlisting}
// Worked example: the exact repro that caught this bug, and what
// Convert::ToInt32(double) returns now that it implements real
// round-half-to-even rounding instead of a bare static_cast<int>().
assert(System::Convert::ToInt32(2.9) == 3);  // was 2 before the fix
assert(System::Convert::ToInt32(3.5) == 4);  // was 3 before the fix -- .NET
                                              // rounds 3.5 up to 4, not to
                                              // the nearer *even* integer,
                                              // since 4 is already even
\end{lstlisting}

\section{Two more findings from the same audit, worth their own worked examples}

Two further findings from that first audit round are worth their own concrete treatment,
because each illustrates a distinct failure shape beyond simple truncation, and this project's
own source shows both bugs and both fixes side by side today.

\cnaclass{Dictionary<TKey, TValue>}'s non-\texttt{const} indexer used to be a plain
\texttt{TValue\&}, which could not distinguish a read access from a write access --- so reading
a missing key silently inserted a default-constructed value, matching
\texttt{std::unordered\_map::operator[]}'s own convention rather than real .NET's, where the
indexer getter unconditionally throws \texttt{KeyNotFoundException} on any missing key. The
quietly damning detail the audit surfaced: this exact bug had already been found and fixed once
in this same codebase, in \cnaclass{Concurrent::ConcurrentDictionary}, via a dedicated
\texttt{ValueProxy} wrapper whose own doc comment states plainly that it does \emph{not} insert
a default --- the fix had simply never been propagated to the plain, non-concurrent
\cnaclass{Dictionary}. It has been now:

\begin{lstlisting}
// Worked example: Dictionary's operator[] today, a ValueProxy return type
// (not a plain TValue&) so a read and a write can be told apart, exactly
// mirroring the precedent ConcurrentDictionary::ValueProxy already set.
[[nodiscard]] ValueProxy operator[](const TKey& key) { return ValueProxy(this, key); }

[[nodiscard]] const TValue& operator[](const TKey& key) const
{
    auto it = map_.find(key);
    if (it == map_.end())
        throw KeyNotFoundException("The given key was not present in the dictionary.");
    return it->second;
}
\end{lstlisting}

\cnaclass{Concurrent::ConcurrentDictionary::GetOrAdd} / \texttt{AddOrUpdate} had a second, more
subtle problem the audit also caught: both held their internal \texttt{std::mutex} for the
\emph{entire} call, including the user-supplied factory callback's own invocation. Since
\texttt{std::mutex} is non-recursive, a factory that reentrantly called back into the very same
dictionary instance --- a realistic memoization pattern, not a contrived edge case --- would
deadlock the calling thread against itself. The fix moves the factory call outside the lock
entirely, matching real .NET's own documented contract for this method precisely (the factory
may run without the lock held and, under real contention, may even be invoked more than once,
with the loser's result discarded):

\begin{lstlisting}
// Worked example: GetOrAdd today -- the factory call is genuinely outside
// both lock_guard scopes, so a reentrant factory can no longer self-deadlock.
TValue GetOrAdd(const TKey& key, std::function<TValue(const TKey&)> factory)
{
    {
        std::lock_guard<std::mutex> lk(mutex_);
        auto it = map_.find(key);
        if (it != map_.end()) return it->second;
    }
    TValue v = factory(key);   // deliberately unlocked -- see the note above
    std::lock_guard<std::mutex> lk(mutex_);
    auto [it, inserted] = map_.emplace(key, std::move(v));
    return it->second;
}
\end{lstlisting}

\section[A second audit, and what sanitizers found that reasoning alone did not]{A second, independent audit, and what sanitizers found that reasoning alone did not}

A further, separate audit produced seven more findings, six fixed in the same session ---
among them a stream-backed ZIP archive implementation that never actually wrote its output
back to the underlying stream in create or update mode, and five of this project's own
destructors (spanning compression streams and XML writing) that could call
\texttt{std::terminate()} if an I/O failure occurred during RAII cleanup, since a
C\texttt{++} destructor that lets an exception escape terminates the program outright rather
than propagating the failure the way a .NET \texttt{Dispose()} call safely could.

A third of that same round's findings deserves its own worked treatment, because it is a real
race between two things that individually look correct: \cnaclass{TaskCompletionSource<TResult>}%
\texttt{::TrySetResult} first atomically claims completion (\texttt{completed\_.%
compare\_exchange\_strong}), \emph{then} calls \texttt{promise\_.set\_value(result)} to actually
publish the result. If \texttt{TResult}'s own copy constructor throws --- a real possibility for
any non-trivial result type, not a contrived one --- that throw happens \emph{after} completion
was already claimed but \emph{before} the underlying \texttt{std::promise}'s shared state ever
becomes ready. Before the fix, nothing caught that exception at this layer, so it propagated
straight out of \texttt{TrySetResult} with the promise left permanently unset: every current and
future waiter on \texttt{GetResult()} or the bridging \cnaclass{Task} blocked forever, and no
later \texttt{TrySetResult} / \texttt{TrySetException} / \texttt{TrySetCanceled} call could ever
re-claim \texttt{completed\_} to unblock them, since it was already (falsely) true.

\begin{lstlisting}
bool TrySetResult(const TResult& result) {
    bool expected = false;
    if (!completed_.compare_exchange_strong(expected, true)) return false;
    try {
        promise_.set_value(result);         // TResult's copy ctor could throw here
    } catch (...) {
        promise_.set_exception(std::current_exception());  // settle the promise anyway
        throw;                              // still propagates to this immediate caller
    }
    return true;
}
\end{lstlisting}

The fix settles the promise with the copy failure itself inside the same \texttt{catch}, so
every waiter still unblocks --- now with an exception to observe rather than a hang --- while
the immediate caller of \texttt{TrySetResult} still sees the same throw it always did, keeping
existing callers that already wrap the call in their own \texttt{try} / \texttt{catch} unaffected.
The general shape worth remembering: an atomic ``claim the transition'' step and the actual
state-publishing step that follows it are two separate operations, and a real exception thrown
strictly between them is exactly the kind of failure window unit tests built around the happy
path will never exercise on their own.

A first-ever full run of the full sanitizer trio --- ThreadSanitizer, AddressSanitizer, and
UndefinedBehaviorSanitizer --- found real bugs no amount of code review alone had caught:
ThreadSanitizer found one genuine production deadlock (a \texttt{Channel} implementation
missing a \texttt{notify\_all} call, a lost-wakeup bug affecting multiple blocked waiters)
alongside three test-only races; AddressSanitizer found one confirmed heap-buffer-overflow (an
aligned-reallocation helper copying the \emph{old}, smaller allocation's size into a newly
enlarged buffer) plus five real memory leaks in XML document node creation; and
UndefinedBehaviorSanitizer found one genuine undefined-behavior call reaching into vendored
\texttt{miniz} code, from an empty vector's \texttt{.data()} pointer reaching a
non-null-parameter-annotated function.

\section{Why this history belongs in a book, not just a changelog}

None of these bugs are individually remarkable --- a truncation-instead-of-rounding bug, a
missed \texttt{notify\_all}, a leaking XML method. What is worth taking from this chapter is
the project's own response to having a fully completed task tracker: rather than reading
``100\% done'' as ``verified correct,'' it treated that state as precisely the moment a fresh,
independently-skeptical audit was most valuable, and built that audit specifically to check
claims against real .NET source rather than against the project's own prior conclusions. This
is the same discipline \texttt{xna4-spec} auditing (Chapter~\ref{ch:xna4-spec-auditing})
brings to CNA's own XNA surface, applied here one layer further down, to the foundation layer
everything else in this book ultimately stands on.
